% Authoritative LaTeX chapter. The selected recipe commands retain their existing IDs and arguments.

\chapter{실전 활용: 핵심 예시 10개}

\label{chap:09_cookbook}

부록~\ref{chap:04_examples}의 활용 예시를 보완하는 핵심 10개를 모았다. 관측기하, 해면 Rayleigh, 에어로졸, 수중 광학, LUT·배치·수렴 확인의 5개 그룹이며 C 호출은 총 13회다. 기존 R 식별자와 출력 파일명을 유지해 실행 기록과 연결한다. 모든 명령은 Quick start 빌드 후 Bash/Linux/WSL의 \texttt{code/\allowbreak{}OCRT\_\allowbreak{}C}에서 실행한다. 양 언어판의 명령과 출력 파일명은 같다.

\noindent\begin{minipage}{\linewidth}
\begin{ManualCode}

export OMP_NUM_THREADS=1
mkdir -p ../../task/manual_cookbook

\end{ManualCode}
\end{minipage}\par

새 결과 폴더와 OCRT 실험·진단 환경변수가 남아 있지 않은 셸을 사용한다. 아래 리다이렉션은 같은 이름의 파일을 덮어쓰므로 수동 재실행 시 폴더명을 바꾼다. 제공한 \texttt{run\_\allowbreak{}recipes.py}는 기존 결과를 거부하고 상속된 \texttt{OCRT\_\allowbreak{}*} 변수를 지운 뒤 사례에 명시한 설정만 켠다. 10개 예시 모두 여섯 기체 총량을 명시적으로 0으로 두며, 해면이 있으면 경계 굴절률은 1.34다. 에어로졸·수중·배치 계산에 필요한 자료를 먼저 설치한다.

\clearpage

\section{핵심 예시 선택 및 실행}

\label{sec:auto:09_cookbook:1}

\label{sec:cookbook_selection}

{\TableFont
\begin{longtable}{@{}>{\raggedright\arraybackslash}p{0.10\linewidth}>{\raggedright\arraybackslash}p{0.54\linewidth}>{\raggedright\arraybackslash}p{0.29\linewidth}@{}}
\toprule
\textbf{ID} & \textbf{주요 목적} & \textbf{비교 기준} \\
\midrule
\endfirsthead
\toprule
\textbf{ID} & \textbf{주요 목적} & \textbf{비교 기준} \\
\midrule
\endhead
\bottomrule
\endfoot
\hyperref[recipe:R01]{R01} & 검은 경계의 분자 산란 기준 & 기준 계산 \\
\hyperref[recipe:R02]{R02} & 거울 방위각과 U 부호 & R01과 비교 \\
\hyperref[recipe:R07]{R07} & 직달 글린트를 분리한 해면 Rayleigh & R01과 비교 \\
\hyperref[recipe:R21]{R21} & 443 nm의 M80C 에어로졸, AOD555 입력 & 자체 AOD 환산 확인 \\
\hyperref[recipe:R26]{R26} & 대기와 순수 해수의 결합 & 수중 기준 계산 \\
\hyperref[recipe:R28]{R28} & CDOM 흡수 추가 & R26과 비교 \\
\hyperref[recipe:R36]{R36} & 전체 a·b·bb 직접 입력 & 독립 계산 \\
\hyperref[recipe:R38]{R38} & 12방향 Rayleigh LUT & R07의 일치 방향 \\
\hyperref[recipe:R39]{R39} & 두 파장 해양 격자와 네이티브 배치 & 내부 3회 호출 비교 \\
\hyperref[recipe:R40]{R40} & 대기·수중 구적의 독립적인 세분화 & R07 및 R26과 비교 \\
\end{longtable}
}

처음에는 R01, R07, R38로 단일 계산과 격자 출력을 확인한다. 아래 도구 명령은 저장소 최상위에서 실행한다. 계획 단계는 솔버를 호출하지 않는다. manifest.json을 확인한 뒤 같은 선택·결과 폴더에 \texttt{-\allowbreak{}-\allowbreak{}run}을 붙인다. 네이티브 Windows에서는 \texttt{python}과 Windows 실행 파일 경로를 사용한다.

\noindent\begin{minipage}{\linewidth}
\begin{ManualCode}

python3 docs/manual/examples/run_recipes.py --list

python3 docs/manual/examples/run_recipes.py \
  --select R01,R02,R07,R38 \
  --ocrt-root code/OCRT_C --exe code/OCRT_C/build/ocrt \
  --outdir task/cookbook_selected --threads 1 --timeout 900

\end{ManualCode}
\end{minipage}\par

사용 가능한 그룹은 geometry, rayleigh, aerosol, ocean, workflow의 5개다. \texttt{-\allowbreak{}-\allowbreak{}select all}은 현재 핵심 10개, 총 13회 C 호출을 선택한다. R39는 내부에서 3회, R40은 2회 호출하며 나머지는 각각 1회다. \texttt{-\allowbreak{}-\allowbreak{}group ocean}처럼 그룹을 선택할 수도 있다. 비교가 필요한 경우 표에 적힌 기준 사례도 함께 선택한다. 도구는 실행 파일·입력 자료 해시, 전체 인자·환경·실행 시간과 stdout/stderr를 기록한다. 사례는 순차 실행하며 \texttt{-\allowbreak{}-\allowbreak{}threads}는 C 프로세스 내부 스레드 수다.

\clearpage

\section{관측기하와 편광}

\label{sec:auto:09_cookbook:2}

\label{sec:cookbook_group:geometry}

\subsection{R01 · 검은 경계의 기준 계산}
\label{sec:auto:09_cookbook:2.1}

\label{recipe:R01}

\textbf{목적. }SZA40\textdegree{}, VZA30\textdegree{}, RAA90\textdegree{}에서 기체 흡수가 없는 분자 산란 기준값을 만든다. 검은 하부 경계는 수면 반사와 수중 광원을 제외한다. 에어로졸 파일과 AOD 옵션은 넣지 않는다.

\noindent\begin{minipage}{\linewidth}
\begin{ManualCode}

./build/ocrt --surface black --sza 40 --wavelength 555 --pressure 1013.25 \
  --vza 30 --raa 90 --gas-column-h2o 0 --gas-column-o3 0 --gas-column-no2 0 \
  --gas-column-o2 0 --gas-column-co2 0 --gas-column-ch4 0 \
  > ../../task/manual_cookbook/R01.txt \
  2> ../../task/manual_cookbook/R01.log

\end{ManualCode}
\end{minipage}\par

\textbf{결과 파일. }\texttt{R01.txt}.

\textbf{확인·해석. }R01.txt의 \texttt{TOA\_\allowbreak{}rho\_\allowbreak{}I/\allowbreak{}Q/\allowbreak{}U}, \texttt{tau\_\allowbreak{}R}, \texttt{orders}, \texttt{conv}를 읽는다. 무차원 TOA 반사도이며 Q나 U가 음수인 것은 정상일 수 있다. \texttt{conv=\allowbreak{}1}은 수렴 플래그이지 정확도의 보증이 아니다. R02는 이 기준에서 RAA만 바꾸어 비교한다.

\clearpage

\subsection{R02 · 거울 방위각과 U의 부호}
\label{sec:auto:09_cookbook:2.2}

\label{recipe:R02}

\textbf{목적. }RAA만 90\textdegree{}에서 270\textdegree{}로 바꾼다. 수평으로 균일하고 방위 방향으로 대칭인 이 장면에서 두 방향은 태양 주평면에 대한 거울상이다.

\noindent\begin{minipage}{\linewidth}
\begin{ManualCode}

./build/ocrt --surface black --sza 40 --wavelength 555 --pressure 1013.25 \
  --vza 30 --raa 270 --gas-column-h2o 0 --gas-column-o3 0 --gas-column-no2 0 \
  --gas-column-o2 0 --gas-column-co2 0 --gas-column-ch4 0 \
  > ../../task/manual_cookbook/R02.txt \
  2> ../../task/manual_cookbook/R02.log

\end{ManualCode}
\end{minipage}\par

\textbf{결과 파일. }\texttt{R02.txt}.

\textbf{확인·해석. }R02.txt와 R01.txt를 비교한다. 수치 오차 범위에서 I와 Q는 같고 U는 부호가 반대여야 한다. U가 0에 가까우면 절대오차를 사용한다. 편광 자료에서 두 방향을 모두 90\textdegree{}로 접으면 U 부호 정보가 사라진다.

\clearpage

\section{해면 Rayleigh 기준}

\label{sec:auto:09_cookbook:3}

\label{sec:cookbook_group:rayleigh}

\subsection{R07 · 거친 검은 해면의 Rayleigh 기준}
\label{sec:auto:09_cookbook:3.1}

\label{recipe:R07}

\textbf{목적. }R01의 검은 경계를 풍속 5 m/s, 굴절률 1.34인 거친 Fresnel 해면과 검은 물로 바꾼다. 직달 선글린트는 분리하지만 다른 수면–대기 상호작용은 남긴다. 이후 Rayleigh 사례의 기준이다.

\noindent\begin{minipage}{\linewidth}
\begin{ManualCode}

./build/ocrt --surface black_fresnel_ocean --sza 40 --wavelength 555 \
  --pressure 1013.25 --vza 30 --raa 90 --wind-speed 5 --n-water 1.34 \
  --gas-column-h2o 0 --gas-column-o3 0 --gas-column-no2 0 --gas-column-o2 0 \
  --gas-column-co2 0 --gas-column-ch4 0 --decouple-sunglint \
  > ../../task/manual_cookbook/R07.txt \
  2> ../../task/manual_cookbook/R07.log

\end{ManualCode}
\end{minipage}\par

\textbf{결과 파일. }\texttt{R07.txt}.

\textbf{확인·해석. }R07.txt는 이 글린트 정의에 따른 해면 Rayleigh 기준값이다. \texttt{TOA\_\allowbreak{}rho\_\allowbreak{}I/\allowbreak{}Q/\allowbreak{}U}를 R01과 비교해 경계 변경의 영향을 본다. 수중 구성성분의 출사 신호는 없으며 결합 해양의 \texttt{Rrs} 산출물이 아니다.

\clearpage

\section{에어로졸과 AOD 기준 파장}

\label{sec:auto:09_cookbook:4}

\label{sec:cookbook_group:aerosol}

\subsection{R21 · 555 nm 기준 AOD의 M80C 에어로졸}
\label{sec:auto:09_cookbook:4.1}

\label{recipe:R21}

\textbf{목적. }443 nm에서 M80C 에어로졸과 AOD555=0.1을 지정한다. 먼저 data release의 \texttt{inputs/\allowbreak{}M80C.mie}를 설치한다. 계산 파장 443 nm와 AOD 기준 파장 555 nm는 서로 다르다. 이 사례는 자체 출력에서 AOD 환산과 결과를 읽는 방법을 보여준다.

\noindent\begin{minipage}{\linewidth}
\begin{ManualCode}

./build/ocrt --surface black_fresnel_ocean --sza 40 --wavelength 443 \
  --pressure 1013.25 --vza 30 --raa 90 --wind-speed 5 --n-water 1.34 \
  --gas-column-h2o 0 --gas-column-o3 0 --gas-column-no2 0 --gas-column-o2 0 \
  --gas-column-co2 0 --gas-column-ch4 0 --decouple-sunglint \
  --mie inputs/M80C.mie --aod-555 0.1 \
  > ../../task/manual_cookbook/R21.txt \
  2> ../../task/manual_cookbook/R21.log

\end{ManualCode}
\end{minipage}\par

\textbf{결과 파일. }\texttt{R21.txt}.

\textbf{확인·해석. }\texttt{AOD\_\allowbreak{}ref=0.1}, \texttt{AOD\_\allowbreak{}ref\_\allowbreak{}nm=555}를 먼저 확인한다. 계산 파장의 \texttt{AOD\_\allowbreak{}band}는 Mie 파일의 소광계수 비로 환산되며 \texttt{AOD\_\allowbreak{}ext\_\allowbreak{}ratio}와 함께 읽는다. 이어서 \texttt{TOA\_\allowbreak{}rho\_\allowbreak{}I/Q/U}와 수렴 표시를 확인한다. R07은 555 nm이고 이 사례는 443 nm이므로 두 반사도의 차이를 에어로졸 기여로 직접 해석하지 않는다. 활성 에어로졸 경로의 수치 기본값도 사용된다.

\clearpage

\section{수중 구성성분과 전체 IOP}

\label{sec:auto:09_cookbook:5}

\label{sec:cookbook_group:ocean}

\subsection{R26 · 분자 대기가 있는 순수 해수}
\label{sec:auto:09_cookbook:5.1}

\label{recipe:R26}

\textbf{목적. }OCRT 구성성분 입력을 모두 명시적으로 0으로 두고 결합 해양을 선택한다. 수온 20\textdegree{}C, 염분 35 g/kg, 경계 굴절률 1.34, 풍속 3 m/s를 고정한다. 1013.25 hPa의 분자는 남고 기체 흡수와 에어로졸은 없다.

\noindent\begin{minipage}{\linewidth}
\begin{ManualCode}

./build/ocrt --surface ocean --sza 40 --wavelength 555 --pressure 1013.25 \
  --vza 30 --raa 90 --wind-speed 3 --n-water 1.34 --gas-column-h2o 0 \
  --gas-column-o3 0 --gas-column-no2 0 --gas-column-o2 0 --gas-column-co2 0 \
  --gas-column-ch4 0 --water-model ocrt --water-temperature 20 \
  --water-salinity 35 --ocrt-chl 0 --ocrt-tsm 0 --ocrt-adom440 0 \
  > ../../task/manual_cookbook/R26.txt \
  2> ../../task/manual_cookbook/R26.log

\end{ManualCode}
\end{minipage}\par

\textbf{결과 파일. }\texttt{R26.txt}.

\textbf{확인·해석. }R26.txt의 \texttt{Rrs0plus\_\allowbreak{}I}, \texttt{rrs0minus\_\allowbreak{}I}, IOP, TOA 성분 키를 읽는다. 단일 텍스트 키는 격자의 \texttt{Rrs\_\allowbreak{}I/\allowbreak{}rrs\_\allowbreak{}I}와 다르다. 이 순수 해수 경로의 내부 태양광 정규화는 pi이므로 비영 구성성분 경로와 비교할 때 원시 Lu 대신 정규화 비율을 사용한다.

\clearpage

\subsection{R28 · CDOM 흡수 추가}
\label{sec:auto:09_cookbook:5.2}

\label{recipe:R28}

\textbf{목적. }R26에 aDOM440=0.1 m\textsuperscript{-1}을 추가하고 Chl과 TSM은 0으로 유지한다. 기본 aDOM 기울기는 0.014 nm\textsuperscript{-1}이며 대기와 해면 조건은 모두 같다.

\noindent\begin{minipage}{\linewidth}
\begin{ManualCode}

./build/ocrt --surface ocean --sza 40 --wavelength 555 --pressure 1013.25 \
  --vza 30 --raa 90 --wind-speed 3 --n-water 1.34 --gas-column-h2o 0 \
  --gas-column-o3 0 --gas-column-no2 0 --gas-column-o2 0 --gas-column-co2 0 \
  --gas-column-ch4 0 --water-model ocrt --water-temperature 20 \
  --water-salinity 35 --ocrt-chl 0 --ocrt-tsm 0 --ocrt-adom440 0.1 \
  > ../../task/manual_cookbook/R28.txt \
  2> ../../task/manual_cookbook/R28.log

\end{ManualCode}
\end{minipage}\par

\textbf{결과 파일. }\texttt{R28.txt}.

\textbf{확인·해석. }\texttt{a\_\allowbreak{}dom}, \texttt{a\_\allowbreak{}total}, 정규화된 \texttt{Rrs0plus\_\allowbreak{}I}를 확인한다. CDOM은 흡수를 추가하며 입자 위상함수를 추가하지 않는다. 비영 구성성분 경로의 내부 태양광 정규화는 1이고 순수 해수 경로는 pi이므로 R26과는 비율·반사도로 비교한다.

\clearpage

\subsection{R36 · 전체 IOP 직접 입력}
\label{sec:auto:09_cookbook:5.3}

\label{recipe:R36}

\textbf{목적. }전체 a=0.1, b=0.2, bb=0.005 m\textsuperscript{-1}을 직접 입력한다. 이미 순수 해수와 모든 구성성분을 포함한 총량이므로 나중에 순수 해수 양을 다시 더하지 않는다.

\noindent\begin{minipage}{\linewidth}
\begin{ManualCode}

./build/ocrt --surface ocean --sza 40 --wavelength 555 --pressure 1013.25 \
  --vza 30 --raa 90 --wind-speed 3 --n-water 1.34 --gas-column-h2o 0 \
  --gas-column-o3 0 --gas-column-no2 0 --gas-column-o2 0 --gas-column-co2 0 \
  --gas-column-ch4 0 --water-model iop --water-temperature 20 \
  --water-salinity 35 --iop-a 0.1 --iop-b 0.2 --iop-bb 0.005 \
  > ../../task/manual_cookbook/R36.txt \
  2> ../../task/manual_cookbook/R36.log

\end{ManualCode}
\end{minipage}\par

\textbf{결과 파일. }\texttt{R36.txt}.

\textbf{확인·해석. }\texttt{a\_\allowbreak{}total}, \texttt{b\_\allowbreak{}total}, \texttt{bb\_\allowbreak{}total}을 확인하고 \texttt{Rrs0plus\_\allowbreak{}I}를 읽는다. 입력은 a>=0, b>0, 0<bb<0.5 b를 만족한다. 외부 위상 파일과 전방 절단 없이 기본 직접 IOP 위상 구성을 사용한다. a/b/bb만으로 가능한 모든 위상함수가 유일하게 정해지는 것은 아니다.

\clearpage

\section{LUT·배치·수렴 확인}

\label{sec:auto:09_cookbook:6}

\label{sec:cookbook_group:workflow}

\subsection{R38 · 12방향 Rayleigh LUT}
\label{sec:auto:09_cookbook:6.1}

\label{recipe:R38}

\textbf{목적. }R07의 단일 방향을 VZA=0/30/60\textdegree{}, RAA=0/90/180/270\textdegree{}로 바꾼다. 한 파장·한 SZA에서 12행 13열인 대기·해면 CSV를 만든다.

\noindent\begin{minipage}{\linewidth}
\begin{ManualCode}

./build/ocrt --surface black_fresnel_ocean --sza 40 --wavelength 555 \
  --pressure 1013.25 --wind-speed 5 --n-water 1.34 --gas-column-h2o 0 \
  --gas-column-o3 0 --gas-column-no2 0 --gas-column-o2 0 --gas-column-co2 0 \
  --gas-column-ch4 0 --decouple-sunglint --lut-vza-step 30 --lut-vza-max 60 \
  --lut-raa-step 90 --output-full-grid ../../task/manual_cookbook/R38.csv \
  > ../../task/manual_cookbook/R38.txt \
  2> ../../task/manual_cookbook/R38.log

\end{ManualCode}
\end{minipage}\par

\textbf{결과 파일. }\texttt{R38.csv}.

\textbf{확인·해석. }실제 격자 헤더 \texttt{rho\_\allowbreak{}I/\allowbreak{}Q/\allowbreak{}U}를 사용한다. VZA30\textdegree{}, RAA90\textdegree{} 행을 R07의 \texttt{TOA\_\allowbreak{}rho\_\allowbreak{}I/\allowbreak{}Q/\allowbreak{}U}와 비교하며 에어로졸 없는 수치 기본값은 같다. 12행과 유한 관측량을 확인하고 RAA270\textdegree{} 및 U 부호를 보존한다. 작은 자료 흐름 시험이며 생산용 각도 해상도 권고는 아니다.

\clearpage

\subsection{R39 · 두 파장 상태: 개별 격자와 네이티브 배치}
\label{sec:auto:09_cookbook:6.2}

\label{recipe:R39}

\textbf{목적. }CDOM 해수의 555와 443 nm 각도 격자를 개별 상태 및 네이티브 해양 배치로 각각 계산한다. 모든 경로에 굴절률 1.34와 수중 구적 96을 명시한다. 명시적 수치 설정을 허용하기 위해서만 ADVANCED를 켠다.

같은 C 작업 폴더에서 네이티브 배치 입력을 먼저 만든다:

\noindent\begin{minipage}{\linewidth}
\begin{ManualCode}

cat > ../../task/manual_cookbook/R39_jobs.csv <<'CSV'
out,wavelength,sza,wind,ocrt_adom440
../../task/manual_cookbook/R39_batch555.csv,555,40,3,0.1
../../task/manual_cookbook/R39_batch443.csv,443,40,3,0.1
CSV

\end{ManualCode}
\end{minipage}\par

\noindent\begin{minipage}{\linewidth}
\begin{ManualCode}

env OCRT_ADVANCED=1 ./build/ocrt --surface ocean --sza 40 --wavelength 555 \
  --pressure 1013.25 --wind-speed 3 --n-water 1.34 --gas-column-h2o 0 \
  --gas-column-o3 0 --gas-column-no2 0 --gas-column-o2 0 --gas-column-co2 0 \
  --gas-column-ch4 0 --water-model ocrt --water-temperature 20 \
  --water-salinity 35 --ocrt-chl 0 --ocrt-tsm 0 --ocrt-adom440 0.1 \
  --n-mu-water 96 --lut-vza-step 30 --lut-vza-max 60 --lut-raa-step 90 \
  --output-full-grid ../../task/manual_cookbook/R39_single555.csv \
  > ../../task/manual_cookbook/R39_single555.txt \
  2> ../../task/manual_cookbook/R39_single555.log

\end{ManualCode}
\end{minipage}\par

\noindent\begin{minipage}{\linewidth}
\begin{ManualCode}

env OCRT_ADVANCED=1 ./build/ocrt --surface ocean --sza 40 --wavelength 443 \
  --pressure 1013.25 --wind-speed 3 --n-water 1.34 --gas-column-h2o 0 \
  --gas-column-o3 0 --gas-column-no2 0 --gas-column-o2 0 --gas-column-co2 0 \
  --gas-column-ch4 0 --water-model ocrt --water-temperature 20 \
  --water-salinity 35 --ocrt-chl 0 --ocrt-tsm 0 --ocrt-adom440 0.1 \
  --n-mu-water 96 --lut-vza-step 30 --lut-vza-max 60 --lut-raa-step 90 \
  --output-full-grid ../../task/manual_cookbook/R39_single443.csv \
  > ../../task/manual_cookbook/R39_single443.txt \
  2> ../../task/manual_cookbook/R39_single443.log

\end{ManualCode}
\end{minipage}\par

\noindent\begin{minipage}{\linewidth}
\begin{ManualCode}

env OCRT_ADVANCED=1 ./build/ocrt --surface ocean --sza 40 --wavelength 555 \
  --pressure 1013.25 --wind-speed 3 --n-water 1.34 --gas-column-h2o 0 \
  --gas-column-o3 0 --gas-column-no2 0 --gas-column-o2 0 --gas-column-co2 0 \
  --gas-column-ch4 0 --water-model ocrt --water-temperature 20 \
  --water-salinity 35 --ocrt-chl 0 --ocrt-tsm 0 --ocrt-adom440 0.1 \
  --n-mu-water 96 --lut-vza-step 30 --lut-vza-max 60 --lut-raa-step 90 \
  --batch-full-grid ../../task/manual_cookbook/R39_jobs.csv \
  > ../../task/manual_cookbook/R39_batch.txt \
  2> ../../task/manual_cookbook/R39_batch.log

\end{ManualCode}
\end{minipage}\par

\textbf{결과 파일. }\texttt{R39\_\allowbreak{}single555.csv}, \texttt{R39\_\allowbreak{}single443.csv}, \texttt{R39\_\allowbreak{}batch555.csv}, \texttt{R39\_\allowbreak{}batch443.csv}.

\textbf{확인·해석. }네 출력 CSV는 각각 12행 60열이다. \texttt{R39\_\allowbreak{}single555}와 \texttt{R39\_\allowbreak{}batch555} 및 443 nm 쌍을 각도·헤더 이름으로 비교한다. 명시적인 n-water는 배치 파장 열에 의한 Quan–Fry 굴절률 자동 선택을 막고 n-mu-water는 배치 48 대 격자 96 기본값 차이를 없앤다. 배치 입력은 BOM 없는 UTF-8이며 경로 안 쉼표와 따옴표를 쓰지 않는다.

네이티브 배치는 CDOM 전용 행에서도 공통 TSM·유기물 자료를 시작 시 미리 읽는다. 일반 단일 Rayleigh 사례와 같은 최소 자료만으로 시작할 수 있다고 가정하지 말고, 설치 장~\ref{chap:02_installation}의 자료 설치 조건을 따른다.

\clearpage

\subsection{R40 · 대기와 수중 구적의 독립적인 수렴 확인}
\label{sec:auto:09_cookbook:6.3}

\label{recipe:R40}

\textbf{목적. }독립적인 수치 세분화를 두 가지 수행한다. R07의 대기 n-mu 24를 32로, 단일 해양 R26의 수중 n-mu-water 48을 96으로 늘린다. 각 쌍은 해당 적분 구적 수만 바꾸며 출력 각도 격자를 바꾸는 실험이 아니다.

\noindent\begin{minipage}{\linewidth}
\begin{ManualCode}

env OCRT_ADVANCED=1 ./build/ocrt --surface black_fresnel_ocean --sza 40 \
  --wavelength 555 --pressure 1013.25 --vza 30 --raa 90 --wind-speed 5 \
  --n-water 1.34 --gas-column-h2o 0 --gas-column-o3 0 --gas-column-no2 0 \
  --gas-column-o2 0 --gas-column-co2 0 --gas-column-ch4 0 --decouple-sunglint \
  --n-mu 32 \
  > ../../task/manual_cookbook/R40_atmos32.txt \
  2> ../../task/manual_cookbook/R40_atmos32.log

\end{ManualCode}
\end{minipage}\par

\noindent\begin{minipage}{\linewidth}
\begin{ManualCode}

env OCRT_ADVANCED=1 ./build/ocrt --surface ocean --sza 40 --wavelength 555 \
  --pressure 1013.25 --vza 30 --raa 90 --wind-speed 3 --n-water 1.34 \
  --gas-column-h2o 0 --gas-column-o3 0 --gas-column-no2 0 --gas-column-o2 0 \
  --gas-column-co2 0 --gas-column-ch4 0 --water-model ocrt \
  --water-temperature 20 --water-salinity 35 --ocrt-chl 0 --ocrt-tsm 0 \
  --ocrt-adom440 0 --n-mu-water 96 \
  > ../../task/manual_cookbook/R40_water96.txt \
  2> ../../task/manual_cookbook/R40_water96.log

\end{ManualCode}
\end{minipage}\par

\textbf{결과 파일. }\texttt{R40\_\allowbreak{}atmos32.txt}, \texttt{R40\_\allowbreak{}water96.txt}.

\textbf{확인·해석. }각 쌍에서 I/Q/U 절대차와 0이 아닌 I의 상대오차를 계산한다. 대기와 해양은 물리 장면이 달라 \texttt{R40\_\allowbreak{}atmos32}와 \texttt{R40\_\allowbreak{}water96}을 직접 비교하지 않는다. 작은 차이는 이 상태에서 이 세분화의 근거이며 보편적인 정확도 한계가 아니다. 필요한 경우 Fourier 차수·층 수·허용치는 따로 세분화한다.

\clearpage

\section{결과의 품질과 재현성 확인}

\label{sec:auto:09_cookbook:7}

\label{sec:cookbook_quality}

도구의 기본 검사는 주요 관측량의 유한수, 예상 격자 행·열·각도, 해양 CSV의 \texttt{water\_\allowbreak{}converged=1}을 확인한다. 결합 해양 단일 텍스트의 \texttt{conv}는 반환된 수중 계산의 수렴을 보고하며 대기·전체 결합의 독립적인 보증은 아니다. 대기 단일 텍스트의 \texttt{conv}는 Fourier SOS 수렴을 모은다. 유효하지 않은 상향 전달비는 따로 기록하고 0으로 바꾸지 않는다. 기본 검사 통과만으로 물리적 정확도가 확립되지는 않는다.

R40은 같은 물리 상태에서 구적 수를 바꾸는 두 비교를 제공한다. 각 비교는 R07 또는 R26을 기준으로 하며 서로 다른 대기·해양 결과를 직접 비교하지 않는다. 실제 적용에는 구적 수뿐 아니라 필요한 층 수·Fourier 차수·수렴 허용치의 안정성과 독립 기준 자료를 확인한다. 보류한 추가 사례의 검토와 검증을 완료하기 전에는 현행 10개에 포함된 것으로 취급하지 않는다.

결과마다 소스·실행 파일·입력 자료 버전, 전체 인자·환경·로그, 파장·기하·기압·풍속·AOD·수중 상태와 단위를 보관한다. 보간은 독립 계산점으로 확인하고, 학습자료는 인접 각도 행만 나누기보다 완전한 대기·해양 상태를 학습·평가 집합으로 나눈다. 실패한 계산을 0으로 채우지 않는다. 이 예시들은 C 실행 파일을 호출하며 Python의 별도 지원 범위와 자료 제한은 부록~\ref{chap:06_python_batch}를 따른다.
